%%% I (Yaniv) left this subsection in the source-material on purpose. We did not feel that it was sufficiently refined to be added to the already very long "documentation of efforts", and it does not improve anything. Yet, let us keep everything together. The curious reader may go over it.

\section{Semi-Definite Programming}
\label{section_Semidefinite_program}

Since the LP was shown to have integrality gap, and it is not clear how to round the variables well, let us consider a semi-definite program (SDP). It will utilize similar intuitions to define the variables and their constraints. Recall that the original intention for defining $Z_{kij}$ was to represent the term $X_{ki} \cdot X_{kj}$ since for $k \in (i \leftrightsquigarrow j)$ we have that $Z_{kij} = 1$ if and only if $X_{ki} = X_{kj} = 1$. Also, since conceptually $X_{ij} \in \{0,1\}$, we can view the single indicators also as multiplications, as $X_{ij} = X_{ij} \cdot X_{ij}$. Then, we can re-define our problem more ``faithfully'', and also add few additional interesting constraints:

\begin{enumerate}
    \item \label{case_SDP_ancestry} We reformulate the ancestry constraints: $\forall i \ne j: X_{ij}^2 + X_{ji}^2 + \sum_{k \in (i \leftrightsquigarrow j)}{X_{ki} \cdot X_{kj}} = 1$. We no longer need the $Z$ variables or their loosely LCA constraints against $X$.
    
    \item \label{case_SDP_nonneg} Non-negativity implies that $X_{ij} \cdot X_{k\ell} \ge 0$ for every choice of $i,j,k,\ell$.

    \item The objective function becomes $\sum_{i}{f_i \cdot D_i} = \sum_{i}{f_i \cdot \sum_{j}{X_{ji}^2}}$.

    \item \label{case_SDP_ijji} We can add a strict constraint $\forall i \ne j: X_{ij} \cdot X_{ji} = 0$ since every STT satisfies it (at most one of the pair is an ancestor of the other). This constraint makes item~\ref{case_SDP_nonneg} unnecessary for $i=k \wedge j=\ell$.

    \item \label{sdp_transitive} By the transitivity of ancestry that holds for every STT, we have that for every triplet $i,j,k$ (all different):
    $X_{ij} \cdot X_{jk} \le X_{ik}$. The right side is not in the form of multiplication, but we can derive from this inequality three loose constraints:
    \begin{enumerate}
        \item $X_{ij} \cdot X_{jk} \le X_{ik} \cdot X_{ik}$: Indicators satisfy $X_{ik} \cdot X_{ik} = X_{ik}$.
        \item $X_{ij} \cdot X_{jk} \le X_{ik} \cdot X_{ij}$: Trivial if $X_{ij}=0$, and if $X_{ij}=1$ then by transitivity.
        \item $X_{ij} \cdot X_{jk} \le X_{ik} \cdot X_{jk}$: Same as the previous case.
    \end{enumerate}

    \item If $X_{ij} = 1$, then every ancestor of $i$ is an ancestor of $j$, but we can also account for $X_{ij}$, and thus $D_i + 1 \le D_j$. If $X_{ij} = 0$ then we cannot lower-bound $D_j$. Combined, we can write $X_{ij} \cdot (1 + D_i) \le D_j$. This seemingly interesting constraint is already taken into consideration as a consequence of the constraints in item-\ref{sdp_transitive}. To see that, By unpacking $D$ back to $X$, and using squaring where needed, we get: $X_{ij}^2 + \sum_{k \in U}{X_{ki} \cdot X_{ij}} \le \sum_{k \in U}{X_{kj}^2}$. For $k \ne i,j$ the inequality holds term-wise. For the remaining terms, we can cancel out $X_{ij}^2$ from both sides, and get: is $X_{ji} \cdot X_{ij}  + X_{ii} \cdot X_{ij} = 0 + X_{ii} \cdot X_{ij} \le X_{jj}^2 = 1$.
    
    \item By ancestry-chain we can also require $X_{ik} \cdot X_{jk} \le X_{ij}^2 + X_{ji}^2$ for \emph{any} triplet of different nodes $i,j,k$. Clarification: For every STT the left side is either $0$ (trivial case) or $1$. If $X_{ik} \cdot X_{jk} = 1$ then both $i$ and $j$ are ancestors of $k$, and one must be the ancestor of the other. The terms on the right side are squared so that we get a proper constraint, with terms of the form $X_{ab} \cdot X_{cd}$.
    
\end{enumerate}

The list above covers all possible choices of triplets $(i,j,k)$: degenerate triplets with only a pair of unique indices ($i,j$) is covered by item \ref{case_SDP_ijji} and the observation that $X_{ij}^2 = X_{ij}$ that we use for other items. There are no variables of the form $X_{ii}$ so there cannot be constraints with a degeneracy from a triplet to a singleton. The rest of the cases cover meaningful relations for non-degenerated triplets. For non-degenerate quadruples $i,j,k,\ell$ there is not much we can require except for non-negativity (item~\ref{case_SDP_nonneg}).

We can translate all of the above to an SDP, where $X_{ij}$ are possibly vectors. The intuition above is for vectors of dimension $1$, but we do not constrain that. We will get an SDP over a matrix of dimension $N \times N$ where $N = \binom{n}{2}$ is the number of pairs, or $X_{ij}$ variables (recall, there are no $X_{ii}$). Due to technical reasons of actual solvers, inequality constraints require us to introduce slack variables, so that instead of (for example) $\vec{a} \cdot \vec{y} \le b$ we will have an equality of the form $\vec{a} \cdot \vec{y} + s = b$. Therefore, the actual size of the matrix would be $N' \times N'$ where $N' = N + S$ and $S$ is the number of inequality constraints. The entries in rows and columns that correspond to slack variables are all zeros except for on the main diagonal. This guarantees non-negative slacks, because the resulting matrix that we solve for is positive semi-definite. See Wikipedia for an example of adding slack variables\footnote{\url{https://en.wikipedia.org/wiki/Semidefinite_programming\#Example_1}}, or our code implementation.

We coded this SDP approach too, with a subset of the constraints that we listed, the ``main'' ones, and got non-integer solutions even for smaller topologies. Even for $n=3$ we start to see non-integer solutions, for example for the weights $f=(4,2,1)$. It is simple to check that the best STT in this case has depths vector $(0,1,2)$. However, the SDP solution yields the depths vector $(0.06,0.94,1.75)$, rounded to two digits. Having non-integer solutions is somewhat expected from a continuous problem, so in order to benefit from the SDP formulation, we need to be able to prove improved guarantees on its rounding.

\subsection{More SDP notes}
One main power of the SDP is our ability to use multiplication terms such as $X_{ab} \cdot X_{cd}$. This is not exactly a simple  multiplication, but if we consider it as such for a moment, then we can define boolean operations on indicator variables quite neatly: $A^2 \equiv AND(X,Y) = X \cdot Y$ and $O^2 \equiv OR(X,Y) = X^2 + Y^2 - X \cdot Y$. Using this, we can also roughly define a multi-input AND, for example: $A_1^2 = AND(a,b)$, $A_2^2 = AND(c,d)$ and $A_3^2 = AND(A_1,A_2)$, so loosely speaking $A_3 = AND(a,b,c,d)$. This is \emph{not} an exact way to represent these logical operations, because there is no commutativity or associativity in our case since we have inner products of vectors that eventually give a scalar result which is the norm of $A_3$ in this example. However, it may be worth to keep this approach in mind, for example to define constraints that effectively ``point out'' the root. Concretely, note that the root $r$ of an STT is the only node that satisfies $AND(X_{r1},\ldots,X_{rn}) = 1$, and also $OR(X_{1r},\ldots,X_{nr}) = 0$ (as discussed in Section~\ref{section_refining_the_LP}, property row-min/column-max).